\section{Related work}
\label{sec:related}

\subsection{Monocular and temporal 3D pose estimation}

Monocular 3D pose estimation must resolve depth from incomplete visual
evidence. Temporal models reduce framewise ambiguity by exploiting motion
context. VideoPose3D demonstrated that dilated temporal convolutions can model
long sequences of 2D keypoints efficiently, including in a semi-supervised
setting \citep{Pavllo2019VideoPose3D}. Temporal attention and multi-scale
dilations have likewise been used to reduce jitter and exploit longer context
\citep{Liu2020TemporalAttention}. Our method also uses dilated temporal
convolutions, but it does not lift a 2D sequence. It receives two existing 3D
streams and learns how to combine their complementary reliability.

SAM 3D Body provides the upstream representation in this study
\citep{Yang2026SAM3DBody}. Its MHR output includes body, hand, and foot joints,
giving a 70-joint sequence after inference. We treat this estimator as fixed.
Consequently, the proposed contribution is neither a new body model nor a new
image-to-pose network; it is a downstream fusion layer whose failure modes and
claims must remain conditional on the upstream estimator.

\subsection{Calibrated multi-view fusion and triangulation}

Classical multi-view pose pipelines recover 3D points from synchronized 2D
correspondences and camera projection matrices. Learnable triangulation adds
confidence weighting or volumetric aggregation to this geometry
\citep{Iskakov2019Learnable}. Cross-view fusion can also be introduced at the
image-feature or 2D-pose stage before a structured 3D recovery step
\citep{Qiu2019CrossView}. Camera-disentangled representations reduce the cost of
volumetric processing while retaining calibrated geometric lifting
\citep{Remelli2020Lightweight}. Generalizable triangulation addresses changes in
camera arrangement \citep{Bartol2022Generalizable}, and probabilistic
triangulation models uncertain camera pose for nominally uncalibrated scenes
\citep{Jiang2023Probabilistic}.

These methods estimate 3D structure from multi-view image evidence. By contrast,
our inputs are already 3D, and the method neither receives projection matrices
nor reprojects to images. The two problem classes are complementary. In this
study, conventional triangulation is retained only as an evaluation
pseudo-reference, not as a target used by the fusion learner.

\subsection{Self-supervised multi-view learning}

Multi-view consistency is a strong source of supervision when 3D labels are
scarce. Cross-View Self-Fusion refines 2D observations across uncalibrated views
and trains a monocular lifting network; only one view is needed at inference
\citep{Kim2022CrossViewSelfFusion}. MetaPose jointly estimates camera and pose
from multiple 2D keypoint sets without 3D supervision
\citep{Usman2022MetaPose}. SelfPose3D uses calibrated multi-view images and 2D
pseudo-poses to learn multi-person 3D localization and pose without manual 2D or
3D labels \citep{Srivastav2024SelfPose3D}.

Our use of ``self-supervised fusion'' is therefore narrower. Two views remain
required at inference, and the output is the final fused 3D keypoint sequence.
No camera is estimated. The supervision comes from controlled perturbation of
the observed 3D streams, high-confidence agreement between them, and explicit
sequence structure. This avoids training on the triangulated evaluation stream,
but it does not make the two monocular inputs statistically independent.

\subsection{Rotation representation and applied motion analysis}

Rotation-aware processing requires care because common low-dimensional rotation
parameterizations contain discontinuities. Continuous neural rotation
representations have been studied explicitly \citep{Zhou2019Rotation}. Rather
than regress an unconstrained rotation parameterization, our network predicts
bounded keypoint residuals and derives a proper $\mathrm{SO}(3)$
pelvis--thorax rotation from the resulting skeleton. Circular differences are
used for the signed axial angle so that the $-\pi/\pi$ boundary does not create a
spurious error.

Applied markerless systems such as OpenCap use calibrated smartphone video,
triangulated 2D keypoints, a Mocap-trained marker model, musculoskeletal
constraints, and laboratory validation \citep{Uhlrich2023OpenCap}. That pipeline
supports biomechanical outputs that are outside the scope of this work. We
study uncalibrated pose-stream fusion and return 3D keypoints only. The method is
intended for robust motion representation and comparative analysis, not for
clinical kinetics or validated joint moments.

\begin{table*}[t]
\centering
\caption{Methodological positioning. ``Two 3D streams'' means that complete
per-view 3D poses, rather than images or 2D keypoints, are the fusion inputs.}
\label{tab:positioning}
\scriptsize
\setlength{\tabcolsep}{3pt}
\begin{tabularx}{\linewidth}{p{0.18\linewidth} p{0.15\linewidth}
  >{\centering\arraybackslash}p{0.11\linewidth}
  >{\centering\arraybackslash}p{0.11\linewidth}
  >{\centering\arraybackslash}p{0.11\linewidth} X}
\toprule
Method family & Fusion input & \makecell{Camera\\parameters} &
\makecell{3D labels\\for training} & \makecell{Two views\\at inference} & Primary output \\
\midrule
Learnable triangulation \citep{Iskakov2019Learnable} & Images/2D heatmaps & Required & Commonly used & Yes & Geometric 3D pose \\
Uncalibrated triangulation \citep{Jiang2023Probabilistic} & Images/2D features & Estimated distribution & Used by embedded estimator & Yes & Camera-aware 3D pose \\
Cross-View Self-Fusion \citep{Kim2022CrossViewSelfFusion} & Multi-view 2D poses & Not required & No & No & Trained monocular estimator \\
MetaPose \citep{Usman2022MetaPose} & Multi-view 2D keypoints & Estimated & No & Yes & Camera and 3D pose \\
SelfPose3D \citep{Srivastav2024SelfPose3D} & Calibrated images and 2D pseudo-poses & Required & No manual labels & Yes & Multi-person 3D pose \\
Proposed method & Two 3D streams & Not used & No & Yes & Fused temporal 3D pose \\
\bottomrule
\end{tabularx}
\end{table*}


Table~\ref{tab:positioning} summarizes the resulting methodological boundary.
The closest prior work shares multi-view self-supervision, but differs in input
space, inference-time view count, and output objective. The novel claim is thus
the combination of post-estimation 3D fusion, explicit rotation preservation,
and label-free temporal training, not any single component such as a temporal
convolution, canonical frame, or skeletal loss.
